\section{Comparison with the complete supplied corpus}
\label{sec:corpus-review}

The comparison is pinned to ProveIt commit
\nolinkurl{17f92105e2bc4317feae6c219cb2143437b482ff}.
At this revision the PA research directory contains the extracted contents
of 36 delivered packets. The incoming directory contains 26 ZIP files:
21 are retained copies of those packets, and five are later cached-controller
packets. Thus there are 41 top-level research packets to compare, rather
than 62 independent constructions. The numerous
\nolinkurl{pa194-*} and \nolinkurl{pa195-*} proof directories develop
different proof layers for the same designated machines; their names do
not signify additional machine designs.

The comparison has two levels. First, a fresh read-only census scans every
\nolinkurl{.tm} and \nolinkurl{.txt} file under the entire PA directory,
every incoming ZIP, and every ZIP recursively encountered in those sources.
Archive payloads are deduplicated by their complete SHA-256. Second, the
reports and source descriptions classify the recognized tables as principal
results, supplementary implementations, search seeds, or rejected
experiments. A small state count or a filename beginning with ``candidate''
is never sufficient for a positive classification.

The census visits 44 distinct ZIP payloads through 71 archive references.
It recognizes 4,417 table occurrences, 620 exact byte identities, and 446
action-preserving, root-free control-graph classes. Two disconnected
synthetic controls have no root reaching their whole graph and are left
unclassified. The numeric tables contain eight 194-state graph classes,
sixteen 195-state classes, and ninety-nine 196-state classes. These are
corpus counts, including diagnostics and repeats of ancestors; they are
not counts of established consistency machines. No recognized complete
table has between 10 and 193 working states. This observation is not a
state-count lower bound.

\subsection{The eight existing 194-state graphs}

\begin{longtable}{@{}p{0.12\linewidth}p{0.20\linewidth}p{0.58\linewidth}@{}}
\toprule Family & Controls / source sites & Change and status\\\midrule\endhead
PA194-67 & 134 + 67 - 7; 3673 sites & Separate diagonal head; 67 AX readers and seven compatible history folds. principal; ordinary PA proof and two symbolic physical checkers\\[4pt]
PA194-AD & 134 + 67 - 7; 3673 sites & Exchange the physical A and D register positions; three changed native rows. principal companion; ordinary transfer and two symbolic physical checkers\\[4pt]
PA194-68 & 134 + 68 - 8; 3673 sites & Different association of the same separate-head emitter, with eight folds. principal; ordinary PA proof with independently reconstructed source/table\\[4pt]
PA194C & 134 + 67 - 7; 3649 sites & Cached Shift/MP controller on PA194-67; six changed native rows. All five later cached packets duplicate this exact dense table. principal; ordinary inherited PA calibration and two symbolic physical checkers\\[4pt]
PA194CS & 135 + 67 - 7 - 1; 3646 sites & Destructive equality operand symmetry shortens Shift by three sites. Extra controller reader 330 is folded with 41; this is not an extra AX fold. principal; source block equivalence and two symbolic physical checkers\\[4pt]
PA194altC & 134 + 68 - 8; 3649 sites & Cached controller composed with the alternate PA194-68 reader allocation. search seed; no separate complete consistency endpoint delivered\\[4pt]
PA194CZ & 134 + 67 - 7; 3649 sites & Replace source-site-26 jump by a zero test of register six; native row (27,1) changes from 0 L 40 to 1 R 21. supplementary seed; primary contracts and positivity-mask invariant, no new independent physical replay\\[4pt]
Rejected scanner194 & rejected & Optimistic shortened-scanner/reader quotient fails the required source simulation. explicitly rejected frame-contract experiment; not a Con(PA) result\\[4pt]
\bottomrule
\end{longtable}


The five principal designs in the table are PA194-67, PA194-AD, PA194-68,
PA194C, and PA194CS. The remaining three are deliberately retained in the
novelty comparison: PA194altC and PA194CZ already occur as supplementary
search seeds, while the shortened-scanner table is explicitly rejected.
All eight classes remain distinct even if one allows tape reflection,
simultaneous read/write symbol complementation, and arbitrary choice of
a graph-reaching root. Symbol complementation is used only as a stronger
test against cosmetic novelty; it does not preserve the specified all-zero
blank input.

Several archive names suggest different cached machines, but their
principal literal result is the same table. Its SHA-256 is
\begin{quote}\small
\nolinkurl{d759313c03b515293fe06b12769fa1771d3829a2bad89febaa9b4391990f3a4b}.
\end{quote}
In particular, recomposing the primary seven-fold decoder with the
Combined controller would reproduce an existing result. The symmetry
of the destructive equality test is also already used by PA194CS.
Likewise, merely installing the Combined controller on the alternate
eight-fold decoder, or adding the zero-register jump at source site 26,
would reproduce supplied search seeds. A proposed fallback must be
compared with these supplementary tables as well as the principal ones.

\subsection{The positive 195-state families and their limitations}

\begin{center}\small
\begin{tabular}{@{}lrrrl@{}}
\toprule
Family & Sites & Outcomes & Startup & Physical organization\\
\midrule
Original PA195 & 3,700 & 4,815 & 58,154 & $134+68-7$\\
195S & 3,673 & 4,780 & 58,154 & $134+67-6$\\
195C & 3,649 & 4,749 & 57,922 & 195S with Combined commands\\
M195S8 & 3,700 & 4,815 & 59,164 & Extra scanner control and eight folds\\
PA195DC & 3,676 & 4,784 & 57,922 & Original diagonal emitter, Combined commands\\
\bottomrule
\end{tabular}
\end{center}

The table identifies the five designated positive manuscript families.
Original PA195 packages the diagonal-vacuity formula with the guarded
gate. The separate-head emitter saves one macro occurrence and 27
primitive sites; it underlies 195S, 195C, and the principal 194-state
families. Combined Shift and MP remove a further 24 primitive sites,
but do not automatically save a reader state. M195S8 uses an additional
scanner control to make a different eight-fold allocation valid.
PA195DC applies the Combined controller to the original bundled-diagonal
source rather than the separate-head source.

The sixteen recognized 195-state graph classes include these five, six
explicitly rejected or diagnostic quotients, a secondary separate195
layout, a crossed allocation, the 195-state parent of PA194CS, and two
secondary implementations in the PA195DC search material. The latter
include a negative-diagonal variant whose author explicitly keeps it
outside the principal consistency theorem. They must not be silently
promoted to independently proved consistency constructions. The full
class inventory records every byte identity and source locator.

\subsection{What the larger designs contribute}

James Chang's PA209 and the Joint/Compatibility U209 ancestor are different
209-state constructions, with respectively 3,948 and 3,909 primitive
source sites. The U208 packets identify two old decoder roles by a
finite-history argument. Entry-context, three-cell, and history
presentations repeatedly describe this same kind of one-state saving;
later helper layouts produce further 208-state graphs. Their equal
state counts do not identify their source programs or transition graphs.

The published PA196 has 3,812 source sites and 4,959 primitive outcomes.
Its allocation comprises 36 tape/register-engine controls, 13 primitive
routers, and 147 readers. The later 196-state proposals change head
ordering, association, physical A/D positions, closing decrements, or
zero-register jump routing. The direct repository catalogue selects
22 rooted 196-state graph classes and explicitly distinguishes two
experimental MP variants and two rejected layouts. The latter both
halt on blank tape after 43,563 transitions; their finite halting traces
are Lean checked in the repository. They are controls, not consistency
machines.

The mathematical source changes show why source length cannot be used
as a proxy for the number of physical states. Diagonal vacuity reduces
the emitter to 3,700 sites, but its first implementation uses 199 states;
later allocations reach 198, 196, and 195. Separating the diagonal head
gives 3,673 sites and eventually reaches 194. An integrated witness gate
uses 3,745 sites and 4,872 primitive outcomes, yet its verified physical
allocations use 206 and then 203 states. Distributed-gate and
short-negation alternatives likewise lose advantageous reader sharing.
The NEG18 separate-emitter source has only 3,441 sites and 4,468
outcomes, but the documented allocation is $134+68-5=197$.
A repaired 157-node representation of the unchanged PA196 source has
a 207-state realization: a smaller abstract DAG can require more
routing and fewer useful contextual identifications.

Changed command protocols form a separate axis. H-only Shift alone has
3,811 sites; Combined196 has 3,788 sites, 4,928 outcomes, and a complete
57,922-step startup. Its successful MP returns the consequent in H,
clears T and D, and does not push a dummy stack value. Shift tests H.
These semantics differ from the original controller, so the command
compiler and arbitrary-incoming-store refutation argument must change
with them. An MP-only bounded register experiment is not the same result
as a complete Combined implementation with all universal source outcomes.

Other packets contain source-only grammar proposals, failed scanner
shortcuts, pruned sources without a calibrated blank seed, and
intermediate parity/control components. Their contribution is evidence
about particular search spaces or reusable proof obligations. A failed
universal frame contract rejects the proposed simulation argument; it
does not by itself prove that the failing frame occurs on the blank run.
Conversely, a long nonhalting prefix does not repair that argument.

\subsection{Proof status and inherited assumptions}

The ordinary PA calibration in the recent reports is stronger than
experimental evidence but distinct from a freshly elaborated proof-assistant
endpoint. It supplies syntactic translations of PA and numerical-calculus
refutations, an explicit command compiler, exhaustive command enumeration,
and positive-duration physical simulations. The reports explicitly address
the code-zero alias $I(0,0)=0$, the Cantor-stack alias $[0;0]=0$, early
halts during proof compilation, an arbitrary carried store between trials,
and the coincident-variable case in the finite bit-comprehension argument.
These points matter to the claimed PA-provable biconditional.

The repository's current README nevertheless states that neither
machine-specific reduction has yet been established in iPRA for the
primary PA194, original PA195, or PA196 candidate. Its ordinary finite-run
Lean theorems and conditional arithmetic endpoints do not close those
obligations. This is consistent with reusing an explicit ordinary
mathematical PA calibration; it must not be described as a completed
Lean/Rocq or iPRA result. The present comparison audits the arguments and
their interfaces but does not independently rebuild every historical
proof-assistant development.

\subsection{Coverage limits and reproducibility}

The literal recognizer requires a UTF-8 text file of five-column rows
\texttt{state read write direction next}, numeric state labels, a binary
alphabet, both rows for every working state, and exactly one undefined
halt target. Blank lines and whole-line comments are allowed. Files using
other syntaxes, multiple undefined halts, code-generated tables not
materialized as text, and tables embedded solely in JSON or Lean are
outside this mechanical recognizer. They are considered through the
reports and source descriptions rather than counted as recognized literal
tables. The archive ledger gives all 41 top-level packet contributions.

Graph classification forgets the designated start and compares each
root that reaches the full working graph by deterministic breadth-first
numbering with symbol-labelled edges and exact write/direction actions.
Taking the least such encoding yields a complete invariant for this
stated class of root-free graphs. This is stronger than comparing native
labels or file hashes, but makes no assertion about semantic equivalence
or correctness. Exact machine bytes, graph hashes, all occurrence paths,
archive identities, the parser, and the read-only reproduction command
are included in the comparison artifacts.

\subsection{Prior work on expanded macro searches}

The supplied search material already contains a partial expansion of
the parameter-copy macros. In
\nolinkurl{pa194_conpa_2026-11-13_2/search/exploratory/expanded_input.txt},
fourteen terminals replace the eleven-terminal basis. Native copy roots
12, 98, and 166 are opened into push root 19 and tail roots 20, 114,
and 179. The old diagonal macro word grows from 191 to 252 terminals.
The associated \nolinkurl{search_expanded.cpp} hashes newly composed
pairs against all fixed pure-reader pairs, so exact fixed-pair reuse
is already present. The recorded expanded and general runs cover
3,363,000 and 1,023,000 iteration slots, respectively, and retain 195
states as their best count. Earlier free-macro-aware grammar searches
also give source-only upper bounds.

Opening deeper interiors of the cons and push macros, using the newer
separate-head or cached source, and performing a physically checked
global optimization can extend those searches. It would be inaccurate
to claim that exposing the three copy tails or recognizing fixed pure
reader pairs is itself a new idea. Nor do those negative runs prove
that a deeper or differently initialized search cannot save a state.
